\chapter{Differential equations}
\label{chap:ode}

Differential equations are where the Banach and the locally convex settings part ways. On a
Banach space the classical theory holds: local solutions exist, are unique, and depend smoothly
(or analytically) on time and initial value. On a Fréchet space even linear equations with
continuous coefficients may have no solution, or more than one
(Chapter~\ref{chap:counterexamples}). What survives is the case in which the state space is a
Banach space and only the \emph{parameters} live in a locally convex space; this is the case
treated in \cite[\S\S2.4--2.5]{GN26}.

We first treat the second order equation of a quadratic spray, which is what the geodesics of
lc-manifolds need, then local flows of vector fields on Banach spaces with $\Ck{n}$, holomorphic
and real analytic dependence, and finally equations with locally convex parameters and their
maximal flows.

\section{Quadratic second order equations}

On a Banach space $E$ we consider the equation of the geodesics of a spray,
\[
  c'' = \Gamma(c)(c', c'),
\]
where $\Gamma \colon U \times E \times E \to E$ is $\Ck{2}$ in the sense of Michal--Bastiani on
an open set $U \subseteq E$ and bilinear in the last two variables. It is equivalent to the first
order system $(c, v)' = (v, \Gamma(c)(v, v))$, whose right hand side is Fr\'echet $\Ck{1}$
(Theorem~\ref{thm:mb-to-frechet}), hence locally Lipschitz.

\begin{theorem}[Local existence, uniform for bounded velocities]
  \label{thm:quadratic-ode-existence}
  \lean{LocallyConvex.Analysis.exists_uniform_isSprayCurveOn_of_norm_le}
  \leanok
  \uses{thm:mb-to-frechet}
  Let $E$ be a Banach space and $x_0 \in U$. There is $r > 0$ such that for every $K \ge 0$ there
  is $\varepsilon > 0$ for which every initial condition $c(0) = y$, $c'(0) = w$ with
  $\|y - x_0\| \le r$ and $\|w\| \le K$ has a solution on $[-\varepsilon, \varepsilon]$ staying in
  $U$.
\end{theorem}

\begin{proof}
  \leanok
  The Picard--Lindel\"of theorem, keeping track of the fact that the solution stays in the ball
  where the vector field is Lipschitz, gives a uniform time for $\|w\| \le \delta$. If $c$ is a
  solution, so is $t \mapsto c(st)$, with initial velocity $s\, c'(0)$; rescaling by
  $s = (K + \delta)/\delta$ gives the statement.
\end{proof}

\begin{theorem}[Uniqueness]
  \label{thm:quadratic-ode-uniqueness}
  \lean{LocallyConvex.Analysis.IsSprayCurveOn.eqOn_Ioo}
  \leanok
  \uses{thm:mb-to-frechet}
  Two solutions on $(a, b)$, staying in $U$, with the same position and velocity at some
  $t_0 \in (a, b)$, agree on $(a, b)$.
\end{theorem}

\begin{proof}
  \leanok
  Integral curves of a locally Lipschitz autonomous field agree near every point where they meet
  (Gr\"onwall); the set where they agree near $t$ is open and closed in $(a, b)$.
\end{proof}

\section{Local flows on Banach spaces}

Let $E$ be a Banach space, $U \subseteq E$ open and $f \colon U \to E$ Fr\'echet $\Ck{n}$,
$1 \le n \le \infty$. We prove that the solutions of $c' = f \circ c$ depend in a $\Ck{n}$ way on
time and initial value, following the implicit function approach of \cite[Ch.~IV, \S1]{Lan99}.

\begin{lemma}[Superposition operators]
  \label{lem:superposition-banach}
  \lean{LocallyConvex.Analysis.contDiffOn_superpose, LocallyConvex.Analysis.hasFDerivAt_superpose}
  \leanok
  Let $K$ be compact. If $f$ is $\Ck{n}$ on $U$, then $\gamma \mapsto f \circ \gamma$ is
  Fr\'echet $\Ck{n}$ on the open set of curves $\gamma \in C(K, E)$ with values in $U$ (sup norm),
  with derivative $\eta \mapsto (s \mapsto f'(\gamma(s))\,\eta(s))$.
\end{lemma}

\begin{proof}
  \leanok
  The derivative comes from the mean value inequality and the uniform continuity of $f'$ near the
  compact set $\gamma(K)$ (compare \cite[Lemma~IV.1.12]{Lan99}). It is again a superposition
  operator, composed with a continuous bilinear map, so induction on $n$ applies.
\end{proof}

\begin{theorem}[Local flows with $\Ck{n}$ dependence]
  \label{thm:banach-local-flow}
  \lean{LocallyConvex.Analysis.exists_contDiffOn_flow, LocallyConvex.Analysis.exists_smoothMBOn_flow}
  \leanok
  \uses{lem:superposition-banach}
  For $x_0 \in U$ there are $\varepsilon > 0$, an open $V \ni x_0$ and
  $\mathrm{Fl} \colon (-\varepsilon, \varepsilon) \times V \to U$, jointly $\Ck{n}$, with
  $\mathrm{Fl}(0, x) = x$ and $\partial_t \mathrm{Fl}(t, x) = f(\mathrm{Fl}(t, x))$. If $f$ is
  smooth in the sense of Michal--Bastiani, so is $\mathrm{Fl}$.
\end{theorem}

\begin{proof}
  \leanok
  Let $J\gamma(s) = \int_0^s \gamma$ on $C([-1, 1], E)$ and
  $\Psi(\lambda, x, \gamma) = (\lambda, x, \gamma - x - \lambda J(f \circ \gamma))$. By
  Lemma~\ref{lem:superposition-banach}, $\Psi$ is $\Ck{n}$, and its derivative at
  $(0, x_0, x_0)$ is a shear, hence invertible. The inverse function theorem gives a $\Ck{n}$
  family $\Gamma(\lambda, x)$ with $\Gamma = x + \lambda J(f \circ \Gamma)$; put
  $\mathrm{Fl}(t, x) = \Gamma(t, x)(1)$. By injectivity of $\Psi$ near $(0, x_0, x_0)$,
  $\Gamma(\sigma\lambda, x)(1) = \Gamma(\lambda, x)(\sigma)$ for $|\sigma| \le 1$, and the
  fundamental theorem of calculus gives the equation. Lang applies the implicit mapping theorem
  only in $x$, on a fixed short time interval \cite[Theorem~IV.1.14]{Lan99}; taking the time
  scale $\lambda$ as a parameter gives joint smoothness directly. We do not claim this argument is
  new.
\end{proof}

\begin{theorem}[Complex time]
  \label{thm:banach-local-flow-rclike}
  \lean{LocallyConvex.Analysis.exists_contDiffOn_flow_rclike}
  \leanok
  \uses{lem:superposition-banach}
  Let $\mathbb{K} \in \{\mathbb{R}, \mathbb{C}\}$, $E$ a Banach space over $\mathbb{K}$ and $f$
  of class $\Ck{n}$ over $\mathbb{K}$ on $U$, $1 \le n \le \infty$. For $x_0 \in U$ there are
  $\varepsilon > 0$, an open $V \ni x_0$ and $\mathrm{Fl} \colon B_{\mathbb{K}}(0, \varepsilon)
  \times V \to U$, jointly $\Ck{n}$ over $\mathbb{K}$, with $\mathrm{Fl}(0, x) = x$, whose
  restriction to real times solves $\partial_t \mathrm{Fl}(t, x) = f(\mathrm{Fl}(t, x))$.
\end{theorem}

\begin{proof}
  \leanok
  The proof of Theorem~\ref{thm:banach-local-flow} works verbatim over $\mathbb{K}$: the curves
  are parametrized by the real interval $[-1, 1]$ and only the time scale $\lambda$ lies in
  $\mathbb{K}$; Mathlib's inverse function theorem holds over $\mathbb{R}$ and $\mathbb{C}$. The
  rescaling identity is used for real $\sigma$ and real times. Theorem~\ref{thm:banach-local-flow}
  is the case $\mathbb{K} = \mathbb{R}$.
\end{proof}

\begin{theorem}[Holomorphic and real analytic flows {\cite[Theorem~2.4.34]{GN26}}]
  \label{thm:banach-local-flow-analytic}
  \lean{LocallyConvex.Analysis.exists_differentiableOn_flow_complex,
    LocallyConvex.Analysis.exists_realAnalyticOn_flow}
  \leanok
  \uses{thm:banach-local-flow-rclike, thm:banach-holomorphic, def:real-analytic,
    lem:normed-complexification, thm:banach-ode-uniqueness}
  If $f$ is complex differentiable on an open subset $U$ of a complex Banach space, its local
  flows are complex differentiable in complex time and initial value. If $f$ is real analytic on
  an open subset $U$ of a real Banach space, its local flows
  $\mathrm{Fl} \colon (-\varepsilon, \varepsilon) \times V \to U$ are real analytic.
\end{theorem}

\begin{proof}
  \leanok
  The holomorphic case is Theorem~\ref{thm:banach-local-flow-rclike} for $\mathbb{K} = \mathbb{C}$,
  since complex differentiable maps are $\Ck{1}$ over $\mathbb{C}$
  (Theorem~\ref{thm:banach-holomorphic}). In the real analytic case, let $g$ be a complex
  analytic extension of $f$ to an open subset of $E_{\mathbb{C}}$, a complex Banach space with
  the norm of Lemma~\ref{lem:normed-complexification}, and $\Phi$ its holomorphic flow. For real
  $t$ and $x \in E$, $t \mapsto \mathrm{Fl}(t, x)$ and $t \mapsto \Phi(t, x)$ solve
  $c' = g \circ c$ with $c(0) = x$, since $g$ extends $f$; by uniqueness they agree. So $\Phi$,
  composed with $(\mathbb{R} \times E)_{\mathbb{C}} \cong \mathbb{C} \times E_{\mathbb{C}}$, is
  a complex analytic extension of $\mathrm{Fl}$. This is the proof of Glöckner and Neeb (real
  analytic case, p.~206), except that they show that complex solutions with real data stay real
  (their Lemma~2.4.30), while we use the real flow and uniqueness.
\end{proof}

\begin{corollary}[Time-dependent equations with parameters {\cite[Theorem~2.4.34(b)]{GN26}}]
  \label{cor:banach-flow-analytic-param}
  \lean{LocallyConvex.Analysis.exists_realAnalyticOn_flow_param}
  \leanok
  \uses{thm:banach-local-flow-analytic}
  Let $E$, $\Lambda$ be real Banach spaces and $f$ real analytic on an open subset $U$ of
  $\mathbb{R} \times E \times \Lambda$. Near every $p_0 \in U$ the solution of
  $c'(t) = f(t, c(t), \lambda)$ with $c(t_0) = x$ is real analytic in $(t, t_0, x, \lambda)$.
\end{corollary}

\begin{proof}
  \leanok
  Apply Theorem~\ref{thm:banach-local-flow-analytic} to the real analytic field
  $(s, x, \lambda) \mapsto (1, f(s, x, \lambda), 0)$. Along its integral curves the first
  component has derivative $1$ and the last one derivative $0$, so the curve through
  $(t_0, x, \lambda)$ is $t \mapsto (t_0 + t, c(t_0 + t), \lambda)$. Glöckner and Neeb allow any
  locally convex parameter space; that case is Theorem~\ref{thm:lc-param-flow}.
\end{proof}

\begin{theorem}[Uniqueness]
  \label{thm:banach-ode-uniqueness}
  \lean{LocallyConvex.Analysis.eqOn_Ioo_of_contDiffOn}
  \leanok
  If $f$ is $\Ck{1}$ on $U$, two solutions of $c' = f \circ c$ on $(a, b)$ with values in $U$
  that agree at one instant agree on $(a, b)$.
\end{theorem}

\begin{proof}
  \leanok
  $f$ is locally Lipschitz; apply the argument of Theorem~\ref{thm:quadratic-ode-uniqueness}.
\end{proof}

\section{Parameters in locally convex spaces and maximal flows}

Outside Banach spaces local existence and uniqueness can fail, even for continuous linear vector
fields \cite[Examples~2.4.2 and 2.4.3]{GN26}.

\begin{lemma}[Uniform contractions with parameters {\cite[Lemma~2.3.26(a), Proposition~2.3.27(a)]{GN26}}]
  \label{lem:uniform-contraction}
  \lean{LocallyConvex.Analysis.UniformContraction,
    LocallyConvex.Analysis.UniformContraction.continuousOn_fixedPoint,
    LocallyConvex.Analysis.UniformContraction.hasDirDerivAt_fixedPoint,
    LocallyConvex.Analysis.UniformContraction.contDiffMBOn_fixedPoint,
    LocallyConvex.Analysis.UniformContraction.dMB_fixedPoint_eq,
    LocallyConvex.Analysis.UniformContraction.contDiffMBOn_fixedPoint_of_contDiffMBOn,
    LocallyConvex.Analysis.UniformContraction.eq_fixedPoint}
  \leanok
  Let $f(p, x)$ be a family of $\theta$-contractions ($\theta < 1$) of a closed set of a Banach
  space $F$, with $p$ in an open subset of a topological vector space. The fixed point
  $\varphi(p)$ is continuous, and if $f$ is $\Ck{1}$ so is $\varphi$, with
  $\dd\varphi(p, q) = \sum_n A_p^n b_p(q)$, $A_p = \dd_2 f(p, \varphi(p))$,
  $b_p = \dd_1 f(p, \varphi(p))$. If $f$ is $\Ck{k}$, so is $\varphi$.
\end{lemma}

\begin{proof}
  \leanok
  Continuity from the uniform contraction estimate; the derivative is the Neumann series, whose
  terms converge locally uniformly (Tannery's theorem). For $\Ck{k}$, by induction on $k$:
  differentiating $\varphi(p) = f(p, \varphi(p))$, $\dd\varphi(p, q)$ is the fixed point of
  the uniform family of $\theta$-contractions $w \mapsto \dd f((p, \varphi(p)), (q, w))$ with
  parameter $(p, q)$, which is $\Ck{k-1}$ by the induction hypothesis applied to $\varphi$.
\end{proof}

\begin{lemma}[$\Ck{k}$ superposition operators with parameters {\cite[Proposition~1.7.5]{GN26}}]
  \label{lem:ck-superposition-param}
  \lean{LocallyConvex.Analysis.pairCLM,
    LocallyConvex.Analysis.contDiffMBOn_superposeParam_of_contDiffMBOn}
  \leanok
  Let $K$ be a compact metric space, $\kappa \colon K \to \mathbb{R}$ continuous, and $H$ a
  $\Ck{k}$ map on an open subset of $\mathbb{R} \times (E' \times Z)$ containing
  $\kappa(K) \times \Omega$, $E'$ normed. The operator
  $(\psi, z) \mapsto (s \mapsto H(\kappa(s), \psi(s), z))$ is $\Ck{k}$ on the open set of
  $(\psi, z) \in C(K, E') \times Z$ with $(\psi(s), z) \in \Omega$ for all $s$.
\end{lemma}

\begin{proof}
  \leanok
  The difference quotients converge uniformly in $s$ by the mean value inequality. By induction on
  $k$: the derivative is the superposition operator of $\dd H(\cdot, (0, \cdot))$, read on
  $C(K, E' \times E') \times (Z \times Z)$.
\end{proof}

\begin{theorem}[Solutions with locally convex parameters {\cite[Theorem~2.4.34]{GN26}}]
  \label{thm:lc-param-flow}
  \lean{LocallyConvex.Analysis.contDiffMBOn_superposeParam,
    LocallyConvex.Analysis.complexAnalyticOn_of_contDiffMBOn_one,
    LocallyConvex.Analysis.ODEFlowParam.complexAnalyticOn_flowPhi,
    LocallyConvex.Analysis.ODEFlowParam.hasDerivAt_flowPhi_time,
    LocallyConvex.Analysis.ODEFlowParam.flowPhi_mem_of_real,
    LocallyConvex.Analysis.exists_complexAnalyticOn_flow_param,
    LocallyConvex.Analysis.exists_realAnalyticOn_flow_lcParam,
    LocallyConvex.Analysis.ODEFlowParam.contDiffMBOn_flowPhi_of_contDiffMBOn,
    LocallyConvex.Analysis.ODEFlowParam.hasDerivAt_flowPhi_time_of,
    LocallyConvex.Analysis.exists_contDiffMBOn_flow_lcParam,
    LocallyConvex.Analysis.ODEFlowParam.exists_bounds_of_lipschitz,
    LocallyConvex.Analysis.ODEFlowParam.hasDerivAt_flowPhi_time_continuous,
    LocallyConvex.Analysis.exists_continuousOn_flow_lcParam}
  \leanok
  \uses{lem:uniform-contraction, lem:ck-superposition-param, def:real-analytic,
    lem:normed-complexification}
  Let $E$ be a complex Banach space, $W$ a complex locally convex space and $f$ complex analytic on
  an open subset of $\mathbb{C} \times E \times W$. The solution $\Phi(t, t_0, y_0, p)$ of
  $y' = f(t, y, p)$, $y(t_0) = y_0$, is complex analytic in all variables near any point. The same
  holds for real analytic $f$ on $\mathbb{R} \times E \times W$, $E$ a real Banach space and $W$ a
  real locally convex space, and if $f$ is $\Ck{k}$, $k \ge 1$, the solution is $\Ck{k}$ in all
  variables; if $f$ is continuous and Lipschitz in $y$, uniformly in the other variables near
  the point, the solution is continuous.
\end{theorem}

\begin{proof}
  \leanok
  As in \cite{GN26}: $\eta(\tau) = \Phi(t_0 + \tau(t - t_0), \ldots) - y_0$ is the fixed point
  of $\eta \mapsto \int_0^{\tau} (t - t_0) f(t_0 + s(t - t_0), \eta(s) + y_0, p)\,ds$ on
  $C([-1, 1], E)$, a uniform family of contractions, $\Ck{1}$ by the superposition results; its
  differential is complex linear, and a real $\Ck{1}$ map with complex linear differentials into a
  Banach space is complex analytic \cite[Theorem~2.1.12]{GN26}. The rescaling identity (2.107)
  gives the differential equation. In the real case, the complexified equation preserves the closed
  set of curves with values in the real subspace, so its fixed point is real for real data; this
  replaces \cite[Lemma~2.4.30]{GN26}. In the $\Ck{k}$ case the map $g$ is $\Ck{k}$ by
  Lemma~\ref{lem:ck-superposition-param}, hence so is its fixed point
  (Lemma~\ref{lem:uniform-contraction}); the construction is written once for $\mathbb{K} \in
  \{\mathbb{R}, \mathbb{C}\}$. For continuous $f$, the bounds come from continuity and the
  Lipschitz condition, and $\Phi(r) = y_0 + \int_{t_0}^r f(v, \Phi(v), p)\,dv$.
\end{proof}

\begin{theorem}[Maximal flows {\cite[Theorem~2.4.29, Lemma~2.5.7, Theorem~2.5.18, Remark~2.4.35]{GN26}}]
  \label{thm:maximal-flow}
  \lean{LocallyConvex.Analysis.MaximalFlow.dom, LocallyConvex.Analysis.MaximalFlow.flow,
    LocallyConvex.Analysis.MaximalFlow.IsSol.eqOn, LocallyConvex.Analysis.MaximalFlow.flow_eq,
    LocallyConvex.Analysis.MaximalFlow.flow_concat, LocallyConvex.Analysis.MaximalFlow.isSol_flow,
    LocallyConvex.Analysis.MaximalFlow.FlowRegularity,
    LocallyConvex.Analysis.MaximalFlow.HasLocalFlows,
    LocallyConvex.Analysis.MaximalFlow.isOpen_dom_and,
    LocallyConvex.Analysis.MaximalFlow.locallyLipschitzY_of_realAnalyticOn,
    LocallyConvex.Analysis.MaximalFlow.isOpen_dom_and_realAnalyticOn_flow,
    LocallyConvex.Analysis.MaximalFlow.flowRegularity_contDiffMBOn,
    LocallyConvex.Analysis.MaximalFlow.locallyLipschitzY_of_contDiffMBOn,
    LocallyConvex.Analysis.MaximalFlow.isOpen_dom_and_contDiffMBOn_flow,
    LocallyConvex.Analysis.MaximalFlow.isOpen_dom_and_smoothMBOn_flow,
    LocallyConvex.Analysis.MaximalFlow.LocallyLipschitzYParam,
    LocallyConvex.Analysis.MaximalFlow.flowRegularity_continuousOn,
    LocallyConvex.Analysis.MaximalFlow.isOpen_dom_and_continuousOn_flow,
    LocallyConvex.Analysis.MaximalFlow.realAnalyticOn_flow_of_forall_isSol}
  \leanok
  \uses{thm:lc-param-flow}
  Let $E$ be a real Banach space, $W$ a topological space and $f$ defined on
  $U \subseteq \mathbb{R} \times E \times W$, locally Lipschitz in $y$. Solutions of
  $y' = f(t, y, p)$ on open intervals are unique and glue to maximal solutions; let
  $\mathrm{Fl}(t, t_0, y_0, p)$ be the maximal flow, on its domain $\Omega$. Then
  $\mathrm{Fl}(t, t_2, \mathrm{Fl}(t_2, t_1, y_1, q), q) = \mathrm{Fl}(t, t_1, y_1, q)$. If the
  equation admits local flows of a regularity class (closed under restriction, gluing and the
  composition below), $\Omega$ is open and $\mathrm{Fl}$ has that regularity on it. In particular,
  if $W$ is a real locally convex space and $f$ is continuous and locally Lipschitz in $y$ (with
  parameters), $\Ck{k}$ ($k \ge 1$), smooth or real analytic on an open set $U$, then $\Omega$ is
  open and $\mathrm{Fl}$ has the same regularity on $\Omega$; if all solutions with initial data in
  $J \times S$ exist on the open interval $J$, $\mathrm{Fl}$ is real analytic on $J \times J \times S$.
\end{theorem}

\begin{proof}
  \leanok
  Uniqueness on open intervals by connectedness, from local uniqueness (Gronwall). Globalization as
  in \cite[Theorem~2.5.18]{GN26}: given $(\tau, t_0, y_0, p) \in \Omega$, the set of times $t_1$
  near which $\Omega$ is a neighbourhood and $\mathrm{Fl}$ is regular contains $t_0$ (a local flow)
  and is open and closed in the interval of the solution: near a limit point $\theta$, with a
  regular $t_2$ close to $\theta$, $\mathrm{Fl}(t, t_1, y_1, q) =
  \Phi_\theta(t, t_2, \mathrm{Fl}(t_2, t_1, y_1, q), q)$ for the local flow $\Phi_\theta$ at
  $\theta$. A real analytic $f$ is $\Ck{1}$ in $(t, y)$ for fixed $p$, hence locally Lipschitz in
  $y$, and admits real analytic local flows by Theorem~\ref{thm:lc-param-flow}; a $\Ck{k}$ map,
  $k \ge 1$, is locally Lipschitz in $y$ by the continuity of $\dd f$, and admits $\Ck{k}$ local
  flows by the same theorem; so does a continuous map with a local Lipschitz condition.
\end{proof}
